\section{System Evaluation and Verification}
\label{sec:evaluation}

To validate the reliability, performance, and correctness of Omni-IPS, we conducted comprehensive evaluations. This includes checking database content consistency, testing REST APIs with Postman, and verifying that the symbolic engines can solve high-school-level proofs in all three domains.

\subsection{Database Synchronization Statistics}
A critical requirement for propositional solving is ensuring that the Neo4j Graph DB rules align exactly with the Qdrant vector spaces. We successfully re-synchronized the entire knowledge base to a pristine state, resulting in a dataset of exactly **127 rules** and **243 facts** across all domains. Table~\ref{tab:db_stats} summarizes this content.

\begin{table}[h]
    \centering
    \caption{Omni-IPS pristine database synchronization statistics.}
    \label{tab:db_stats}
    \begin{tabular}{lccc}
        \toprule
        \textbf{Domain} & \textbf{Rules (Neo4j \& Qdrant)} & \textbf{Facts (Qdrant Payloads)} & \textbf{Total Entities} \\
        \midrule
        Chemistry & 42 & 84 & 126 \\
        Geometry & 45 & 92 & 137 \\
        Algebra & 40 & 67 & 107 \\
        \midrule
        \textbf{Total} & \textbf{127} & \textbf{243} & \textbf{370} \\
        \bottomrule
    \end{tabular}
\end{table}

This synchronized architecture is indexed correctly in Qdrant, as shown in Figure~\ref{fig:qdrant_collections} (sourced from \texttt{assets/qdrant\_collections.png}).

\begin{figure}[h]
    \centering
    \includegraphics[width=0.7\textwidth]{qdrant_collections.png}
    \caption{Active Qdrant collections for Chemistry, Geometry, and Algebra, confirming the exact ingestion schema.}
    \label{fig:qdrant_collections}
\end{figure}

\subsection{Postman API Verification Template}
We validated backend integrations by sending REST payloads to the API gateway. Listing~\ref{lst:postman_req} details a sample JSON request sent to \texttt{/api/solve} to prove the Pythagorean converse.

\begin{lstlisting}[style=jsonstyle, label={lst:postman_req}, caption={Postman request JSON payload sent to the API gateway at \texttt{/api/solve}.}]
{
  "query": "Prove Triangle(A,B,C) is a Right Triangle given that BC^2 = AB^2 + AC^2",
  "domain": "geometry"
}
\end{lstlisting}

The gateway returned the successful proof validation trace in Listing~\ref{lst:postman_res}. Notice that the goal is solved in exactly **1 logical step** by matching the converse Pythagorean theorem.

\begin{lstlisting}[style=jsonstyle, label={lst:postman_res}, caption={Postman response JSON returned by the backend, detailing the exact rule activation trace.}]
{
  "query": "Prove Triangle(A,B,C) is a Right Triangle given that BC^2 = AB^2 + AC^2",
  "domain": "geometry",
  "mapped_initial_facts": [
    "Triangle(A,B,C)",
    "BC^2=AB^2+AC^2"
  ],
  "mapped_goal": "RightTriangle(A,B,C)",
  "goal_reached": true,
  "applied_rule_ids": [
    "geo_pythagoras_converse"
  ],
  "execution_trace": [
    {
      "rule_id": "geo_pythagoras_converse",
      "fired_rule_repr": "geo_pythagoras_converse [Converse of Pythagorean Theorem]: Triangle(A,B,C) + BC^2=AB^2+AC^2 -> RightTriangle(A,B,C) + RightAngle(Angle(BAC))",
      "new_facts": [
        "RightTriangle(A,B,C)",
        "RightAngle(Angle(BAC))"
      ]
    }
  ],
  "known_facts": [
    "BC^2=AB^2+AC^2",
    "RightAngle(Angle(BAC))",
    "RightTriangle(A,B,C)",
    "Triangle(A,B,C)"
  ]
}
\end{lstlisting}

\subsection{Multi-Domain Proof Traces}
We verified system behavior across multiple complex reasoning scenarios:

\subsubsection{Chemistry Scenario: Preciptation of Barium Sulfate}
*   \textbf{Query}: ``Deduce the precipitate when Barium Chloride reacts with Sodium Sulfate.''
*   \textbf{Initial Mapped Facts}: $\text{BaCl}_2, \text{Na}_2\text{SO}_4$.
*   \textbf{Goal Fact}: $\text{BaSO}_4$.
*   \textbf{Solver Execution}: The forward chainer fires the double displacement rule:
    \begin{equation}
        \text{rxn\_double\_bacl2\_na2so4}: \text{BaCl}_2 + \text{Na}_2\text{SO}_4 \to \text{BaSO}_4 \downarrow + 2\text{NaCl}
    \end{equation}
    which yields $\text{BaSO}_4$ and $\text{NaCl}$ in working memory, successfully satisfying the goal. The stream output for this chemical proof is shown in Figure~\ref{fig:chem_baso4} (derived from \texttt{assets/demo\_BaSO4\_1.png}).

\begin{figure}[h]
    \centering
    \includegraphics[width=0.85\textwidth]{demo_BaSO4_1.png}
    \caption{Interactive stream visualization verifying the precipitation of Barium Sulfate ($\text{BaSO}_4$) from $\text{BaCl}_2$ and $\text{Na}_2\text{SO}_4$.}
    \label{fig:chem_baso4}
\end{figure}

\subsubsection{Algebra Scenario: Equality Substitions}
*   \textbf{Query}: ``Solve x + 2 = 5.''
*   \textbf{Initial Mapped Facts}: $x+2=5, \text{Subtract}(2, \text{both\_sides})$.
*   \textbf{Goal Fact}: $x=3$.
*   \textbf{Solver Execution}: The forward chainer matches the linear equation rule:
    \begin{equation}
        \text{alg\_sub\_two}: \text{Equation}(x+2, 5) \land \text{Subtract}(2) \implies \text{Equation}(x, 3)
    \end{equation}
    proving $x=3$ in a single logical transition with 100\% mathematical validity.
